% Section 14: 01-vector-derivative | Part: 2 | Depends on: see docs/REWRITE_PLAN.md | Status: todo | Original lines: 2301-2335
%!TEX root = ../../../main.tex
\begin{definition}[Vector derivative]
The \emph{vector derivative} (Dirac operator) is $\vd=e^i\partial_i$. Its action on a multivector field $F$ is its geometric product with $F$:
\[
 \vd F=e^i\,\partial_iF=\vd\lc F+\vd\wedge F,\qquad \vd\lc F:=e^i\lc\partial_iF,\quad \vd\wedge F:=e^i\wedge\partial_iF .
\]
\end{definition}

For a $k$-vector field $F$, $\vd\wedge F$ has grade $k+1$ (``curl-like'') and $\vd\lc F$ has grade $k-1$ (``divergence-like''); this is formula \eqref{eq:aA} with the vector $a$ replaced by the operator $\vd$.

\begin{theorem}[$\dd$ is $\vd\wedge$]\label{thm:dGA}
Under the identification above:
\begin{enumerate}
  \item $\dd\omega\ \longleftrightarrow\ \vd\wedge F$.
  \item (Euclidean case) the codifferential $\codiff=(-1)^{n(k+1)+1}\hodge\dd\hodge$ corresponds to $-\vd\lc F$; hence $\vd=\dd-\codiff$ acting on forms.
  \item $\vd\wedge(\vd\wedge F)=0$ and $\vd\lc(\vd\lc F)=0$.
  \item $\vd^2=\vd\cdot\vd=\sum_i\varepsilon_i\partial_i^2$ (the Laplacian, or d'Alembertian), and $\vd^2=\vd\lc(\vd\wedge\,\cdot\,)+\vd\wedge(\vd\lc\,\cdot\,)=-(\dd\codiff+\codiff\dd)$ in the Euclidean case.
  \item (Duality) in terms of the dual of Definition~\ref{def:dual}:
  \[
   \dual{(\vd\wedge F)}=\vd\lc\dual F,\qquad \dual{(\vd\lc F)}=\vd\wedge\dual F .
  \]
\end{enumerate}
\end{theorem}
\begin{proof}
(1) The frame is constant, so $\partial_i$ commutes with it and
$\vd\wedge F=e^j\wedge\partial_jF=\sum_{I,j}\partial_j\omega_I\,e^j\wedge e^{i_1}\wedge\cdots\wedge e^{i_k}$, which is the image of \eqref{eq:dcoord}.
(2) The standard formula for the codifferential is $\codiff=-\sum_i\iota_{e_i}\partial_{i}$ in Euclidean space, and $\iota_{e_i}=e_i\lc$ on multivectors; so $\codiff=-\vd\lc$. Then $\vd=\vd\lc+\vd\wedge=-\codiff+\dd$.
(3) $\vd\wedge(\vd\wedge F)=e^i\wedge e^j\wedge\partial_i\partial_jF=0$ since $e^i\wedge e^j$ is antisymmetric and $\partial_i\partial_j$ symmetric (this is $\dd^2=0$ again). Similarly, by Proposition~\ref{prop:contr}(3), $\vd\lc(\vd\lc F)=(e^i\wedge e^j)\lc\partial_i\partial_jF=0$.
(4) $\vd\vd=e^ie^j\partial_i\partial_j=(e^i\cdot e^j)\partial_i\partial_j$ because the antisymmetric part $e^i\wedge e^j\partial_i\partial_j$ cancels. Expanding $\vd(\vd F)$ by grades and using (3) leaves the two cross terms, which equal $-(\dd\codiff+\codiff\dd)$ by (1)--(2).
(5) is Theorem~\ref{thm:dual}(3) applied with $A=\vd$ (the identities are linear in $A$ and the derivatives $\partial_i$ commute with the constant dual $\pss^{-1}$).
\end{proof}

\begin{remark}
Parts (2) and (4) say that the Dirac operator $\vd$ is a ``square root'' of the Laplacian, $\vd^2=-\Delta_{\mathrm{Hodge}}$, which splits as $\vd=\dd-\codiff$. Item (5) shows why the codifferential is $\pm\hodge\dd\hodge$: it is the \emph{dual} of the exterior derivative, just as contraction is the dual of the wedge.
\end{remark}
